% Part: computability
% Chapter: computability-theory
% Section: rice-theorem

\documentclass[../../../include/open-logic-section]{subfiles}

\begin{document}

\olfileid{cmp}{thy}{rce}
\olsection{రైస్ సిద్ధాంతం}

ఆలోచిస్తే, \olref[tot]{prop:total} నిరూపణలో $\fn{Tot}$కు ప్రత్యేకమైన
వివరాలు ఏ పాత్రా పోషించలేదని తెలుస్తుంది. $x$ అనేది $K$లో ఉన్నప్పుడు
మాత్రమే స్థిర ప్రమేయం $j(y)=0$లాగా పనిచేసేలా $h(x,y)$ను రూపొందించాం;
కాని ఆ పరిస్థితుల్లో దాన్ని మరే పాక్షిక గణనీయ ప్రమేయంలాగైనా పనిచేసేలా
చేయవచ్చు. ఈ పరిశీలన, స్థూలంగా చెప్పాలంటే, గణనీయ ప్రమేయాల ఏ అల్పేతర
ధర్మమూ నిర్ణయించదగినది కాదని చెప్పే మరింత సాధారణ సిద్ధాంతాన్ని
ప్రకటించడానికి వీలు కల్పిస్తుంది.

$\cfind{0}$, $\cfind{1}$, $\cfind{2}$,~\dots{} అనేది పాక్షిక గణనీయ
ప్రమేయాలకు మన ప్రామాణిక లెక్కింపు అని గుర్తుంచుకోండి.

\begin{thm}[రైస్ సిద్ధాంతం]
  $C$ అనేది పాక్షిక గణనీయ ప్రమేయాల ఏ సమితైనా అనుకుందాం; అలాగే
  $A=\Setabs{n}{\cfind{n}\in C}$ అని ఉంచండి. $A$ గణనీయమైతే, $C$
  శూన్య సమితి కావాలి లేదా పాక్షిక గణనీయ ప్రమేయాలన్నిటి సమితి కావాలి.
\end{thm}

$n$, $m$ అనే సూచికలు ఒకే ప్రమేయాన్ని ``గణిస్తే,'' $n$, $m$ రెండూ
$A$లో ఉండాలి లేదా రెండూ దానిలో ఉండకూడదు అనే ధర్మం గల సమితి~$A$ను
\emph{సూచిక సమితి} అంటాం. సిద్ధాంతంలోని సమితి~$A$కు ఈ ధర్మం ఉందని
చూడటం కష్టం కాదు. విలోమంగా, $A$ ఒక సూచిక సమితి, $C$ అనేది ఆ
సూచికలు గణించే ప్రమేయాల సమితి అయితే $A=\Setabs{n}{\cfind{n}\in C}$.

\begin{explain}
ఈ పరిభాషలో, ఏ అల్పేతర సూచిక సమితీ నిర్ణయించదగినది కాదనే ప్రకటనకు
రైస్ సిద్ధాంతం తుల్యం. సిద్ధాంతం ఏమి చెబుతుందో అర్థం చేసుకోవడానికి,
\emph{ప్రోగ్రాములు} (ఉదాహరణకు మీకు నచ్చిన ప్రోగ్రామింగ్ భాషలోని
ప్రోగ్రాములు) మరియు అవి గణించే ప్రమేయాల మధ్య తేడాను నొక్కిచెప్పడం
ఉపయోగకరం. వాక్యనిర్మాణ వస్తువులైన ప్రోగ్రాముల (సూచికల) గురించి కొన్ని
ప్రశ్నలు తప్పకుండా గణనీయాలే: ఈ ప్రోగ్రాములో 150కంటే ఎక్కువ సంకేతాలు
ఉన్నాయా? 22కంటే ఎక్కువ పంక్తులు ఉన్నాయా? ``while'' ప్రకటన ఉందా?
``print'' ప్రకటనకు ఆర్గుమెంటుగా ``hello world'' అనే అక్షరమాల ఎప్పుడైనా
కనిపిస్తుందా? ప్రోగ్రాము యొక్క \emph{ప్రవర్తన} గురించి ఏ అల్పేతర
ప్రశ్నా గణనీయం కాదని రైస్ సిద్ధాంతం చెబుతుంది. కింది ప్రశ్నలు కూడా
ఇందులో చేరతాయి: ప్రోగ్రాము ఇన్‌పుట్~$0$పై ఆగుతుందా? అది ఎప్పుడైనా
ఆగుతుందా? అది ఎప్పుడైనా సరిసంఖ్యను అవుట్‌పుట్ చేస్తుందా?
\end{explain}

\begin{proof}[రైస్ సిద్ధాంత నిరూపణ]
$C$ శూన్య సమితి కాదనీ, పాక్షిక గణనీయ ప్రమేయాలన్నిటి సమితి కూడా
కాదనీ అనుకుందాం; $A$ అనేది~$C$లోని ప్రమేయాల సూచికల సమితి అనుకుందాం.
$A$ గణనీయమైతే ఆగే సమస్యను పరిష్కరించగలమని చూపుతాం; కాబట్టి $A$
గణనీయం కాదు.

సామాన్యతకు నష్టం లేకుండా, ఎక్కడా నిర్వచితం కాని ప్రమేయం~$f$ అనేది
$C$లో లేదని ఊహించవచ్చు (అది $C$లో ఉంటే కింది వాదనలో $C$ను దాని
పూరకంతో మార్చండి). $C$లోని ఏదైనా ప్రమేయం~$g$ను తీసుకోండి. $A$ను
నిర్ణయించగలిగితే, $f$ను గణించే సూచికలకూ $g$ను గణించే సూచికలకూ
మధ్య తేడా చెప్పగలం; అప్పుడు ఆ సామర్థ్యంతో ఆగే సమస్యను పరిష్కరించగలం
అన్నదే భావన.

ఇదిగో విధానం. సార్వత్రిక పాక్షిక గణనీయ ప్రమేయాలను వాడి, కింది
ప్రమేయాన్ని నిర్వచించవచ్చు:
\[
h(x,y) \simeq
\begin{cases}
\text{అనిర్వచితం} & \text{$\cfind{x}(x) \fundefined$ అయితే} \\
g(y) & \text{ఇతర సందర్భాల్లో.}
\end{cases}
\]
$h$ను గణించడానికి ముందుగా $\cfind{x}(x)$ను గణించడానికి ప్రయత్నిస్తాం;
ఆ గణన ఆగితే $g(y)$ను గణించడానికి కొనసాగుతాం; \emph{ఆ} గణన కూడా
ఆగితే దాని అవుట్‌పుట్‌ను తిరిగి ఇస్తాం. మరింత కచ్చితంగా,
\[
h(x,y) \simeq \Proj{2}{0}(g(y),\fn{Un}(x,x)).
\]
అని రాయవచ్చు; ఇక్కడ $\Proj{2}{0}(z_0,z_1)=z_0$ అనేది $0$వ
ఆర్గుమెంటును తిరిగి ఇచ్చే గణనీయమైన $2$-స్థాన ప్రక్షేపణ ప్రమేయం.

అప్పుడు $h$ పాక్షిక గణనీయ ప్రమేయాల సంయుక్తం; కుడివైపు
$\fn{Un}(x,x)$, $g(y)$ రెండూ నిర్వచితమైనప్పుడు మాత్రమే నిర్వచితమై,
$g(y)$కు సమానమవుతుంది.

స్థిరమైన~$x$కు $\cfind{x}(x)$ అనిర్వచితమైతే ప్రతి~$y$కు $h(x,y)$
అనిర్వచితం; $\cfind{x}(x)$ నిర్వచితమైతే $h(x,y)\simeq g(y)$ అని
గమనించండి. అందువల్ల $x$ యొక్క ప్రతి స్థిర విలువకు $h(x,y)$ అనేది
$f$లాగా లేదా $g$లాగా పనిచేస్తుంది. ఈ రెండు సందర్భాల్లో ఏది
నెరవేరుతుందో నిర్ణయించడమే $\cfind{x}(x)$ నిర్వచితమో కాదో నిర్ణయించడం.
కాని ఇది $h_x(y)\simeq h(x,y)$ అనేది~$C$లో ఉందో లేదో నిర్ణయించడమే;
$A$ గణనీయమైతే మనం కచ్చితంగా అదే చేయగలం.

మరింత కచ్చితంగా, $h$ పాక్షిక గణనీయం కనుక ఏదో సూచిక~$e$కు అది
$\cfind{e}$ ప్రమేయమే. $s$-$m$-$n$ సిద్ధాంతం ప్రకారం, ప్రతి~$x$కు
$\cfind{s(e,x)}(y)=h_x(y)$ అయ్యే ఆదిమ పునరావృత్త ప్రమేయం~$s$ ఒకటి
ఉంది. ఇప్పుడు ప్రతి~$x$కు, $\cfind{x}(x)\fdefined$ అయితే
$\cfind{s(e,x)}$ అనేది~$g$తో ఒకే ప్రమేయం; కనుక $s(e,x)$ అనేది~$A$లో
ఉంటుంది. మరోవైపు $\cfind{x}(x)\uparrow$ అయితే $\cfind{s(e,x)}$
అనేది~$f$తో ఒకే ప్రమేయం; కనుక $s(e,x)$ అనేది~$A$లో ఉండదు. మరో
మాటలో, ప్రతి~$x$కు $x\in K$ అయితేనే $s(e,x)\in A$. $A$ గణనీయమైతే
$K$ కూడా గణనీయం అవుతుంది; ఇది వైరుధ్యం. కనుక $A$ గణనీయం కాదు.
\end{proof}

రైస్ సిద్ధాంతం చాలా శక్తివంతమైనది. దాని కొన్ని నమూనా అనువర్తనాలను
కింది తక్షణ ఉపసిద్ధాంతం చూపుతుంది.
\begin{cor}
కింది సమితులు అనిర్ణయనీయాలు.
\begin{enumerate}
\item $\Setabs{x}{\text{$17$ అనేది $\cfind{x}$ వ్యాప్తిలో ఉంది}}$
\item $\Setabs{x}{\text{$\cfind{x}$ స్థిర ప్రమేయం}}$
\item $\Setabs{x}{\text{$\cfind{x}$ సర్వనిర్వచితం}}$
\item $\Setabs{x}{\text{$y<y'$ మరియు $\cfind{x}(y)\fdefined$,
    $\cfind{x}(y')\fdefined$ రెండూ అయినప్పుడల్లా,
    $\cfind{x}(y)<\cfind{x}(y')$}}$
\end{enumerate}
\sourcecorrection{OLTECOMTHY-017}{నాలుగవ సమితి గురించిన ఆంగ్ల మూల వాక్యం ``whenever''తో మొదలైన షరతులో $\cfind{x}(y)\fdefined$ను వేరుగా సంయోజించి, తరువాత $\cfind{x}(y')\fdefined$కు మాత్రమే ``if''ను పెట్టింది. ఉద్దేశించిన అల్పేతర సూచిక ధర్మం---రెండు విలువలూ నిర్వచితమైనప్పుడు నిర్వచన క్షేత్రంపై కఠిన వృద్ధి---స్పష్టమయ్యేలా షరతును సరిచేశాం.}
\end{cor}

\begin{proof}
  ఇవన్నీ అల్పేతర సూచిక సమితులు.
\end{proof}

\end{document}
